\clearpage\thispagestyle{empty}
\pdfbookmark[0]{Quellen zu Komet (24. Mai 1819)}{cometmay1819-sources}
\begin{center}\textbf{Quellen zu Komet (24. Mai 1819)}\end{center}
\noindent Diese KI-Übertragung ins heutige Deutsch folgt \textit{Comet} in den \href{https://gdz.sub.uni-goettingen.de/id/PPN236007467}{\textit{Werken}, Band~VI} (Göttingen, 1874), S.~420--422. Die Werkausgabe nennt als Quelle die \textit{Göttingischen gelehrten Anzeigen}, Stück~83, 24. Mai 1819, S.~825--829. Der Text setzt die Darstellung im \hyperlink{native-422}{Bericht vom 18. Februar 1819} fort. Er verwendet die erste Person Plural und spricht von Gauß in der dritten Person.
\par\smallskip\noindent Die Übertragung übernimmt die ungewöhnliche Formulierung „des einen dieses Kometen“ am Anfang, den Punkt vor dem mit „dass“ beginnenden Satzteil, den Doppelpunkt nach 53 in der zweiten Tabelle, die wiederholte Zeiteinheit in der Seeberger Zeit und die Zahl 32 ohne Strich in der ersten aus den neueren Elementen berechneten Position. KI-Übertragung und KI-Satz.
\clearpage
